\begin{solapartat}
L'alumne pot utilitzar qualsevol criteri de signe coherent (Aquí utilitzem el criteri DIN):

\punts{0,25} L'objecte, llibre, està a 35 cm davant de la lent (s=-35cm),

\punts{0,25} La imatge ha d'estar a 70cm també davant de la lent (s'=-70cm) per poder-la enfocar.

\punts{0,25} Per tant, utilitzant l'equació de la lent prima $\frac{1}{s'} - \frac{1}{s} = \frac{1}{f}$, tenim:
\[ \frac{1}{-70} - \frac{1}{-35} = -\frac{1}{70} + \frac{2}{70} = \frac{1}{70} = \frac{1}{f'}\,. \]
Per tant $f' = 70\,\mathrm{cm}$

\punts{0,25} La potència és la inversa de la distància focal:

\punts{0,25} $P = \frac{1}{f'} = \frac{1}{0,7} = 1,43\ D\ (\text{diòptries})$
\end{solapartat}

\begin{solapartat}
\punts{0,25} L'objecte està a 25 cm davant de la lent (s=-25cm) i la focal de la lent és 20cm. Per tant, utilitzant $\frac{1}{s'} - \frac{1}{s} = \frac{1}{f'}$, tenim:
\[ \frac{1}{s'} - \frac{1}{-25} = \frac{1}{20} \quad\text{i}\quad \frac{1}{s'} = \frac{1}{20} - \frac{1}{25} = \frac{5}{100} - \frac{4}{100} = \frac{1}{100}\,. \]
Per tant, $s' = 100\ \mathrm{cm}$

\punts{0,25} La mida de la imatge la podem trobar amb la relació $\frac{y'}{y} = \frac{s'}{s}$, i, per tant,
\[ y' = y\,\frac{s'}{s} = 10\cdot\frac{100}{-25} = -40\ cm. \]

\punts{0,25} +gran/real/invertida

\punts{0,5} diagrama de rajos: Utilitzem rajos que surten d'un punt de l'objecte per trobar la imatge. El raig paral·lel a l'eix és desviat per la lent i passa per $f'$. El raig que passa pel centre de la lent no es desvia. La imatge es crea on es creuen els dos rajos.
\figura[0.55]{sol_fig1.png}
\end{solapartat}
